% Contenido condensado de cinco páginas para la entrega en Moodle (tarea 1754, Actividad 4).
\fontsize{12}{14.4}\selectfont\setstretch{1.5}\justifying
\begingroup
\noindent\textbf{Resumen}\par
Se distinguen parámetro, estimador y estimación, se explican insesgamiento, eficiencia y consistencia, y se comparan estimaciones puntuales e intervalares en la administración. Tres fichas analizan las propiedades en compras, los usos en áreas funcionales y el intervalo de ventas semanales entre \$18,000 y \$22,000 al 95\% de confianza.\footnote{\textbf{Declaración de uso de IA.} \estadisticadeclaracionia}
\par\smallskip
\renewcommand{\contentsname}{Tabla de contenido}
\setcounter{tocdepth}{1}
\setlength{\cftbeforetoctitleskip}{2pt}\setlength{\cftaftertoctitleskip}{2pt}
\setlength{\cftbeforesecskip}{1pt}
\renewcommand{\cfttoctitlefont}{\normalfont\fontsize{12}{14.4}\selectfont\bfseries\color{itescaRedDark}}
\tableofcontents
\endgroup
\smallskip
\section{Introducción}
Planear compras, contratar personal o asignar un presupuesto exige decidir sin observar a toda la población. La estimación estadística aproxima parámetros mediante muestras y cuantifica su incertidumbre. Su utilidad administrativa depende tanto de las propiedades del estimador como de la representatividad de los datos: una cifra precisa puede orientar mal una decisión si el marco excluye sistemáticamente a parte del mercado \citep{holmesBusiness2023,bradley2021}.

\section{De los estimadores a las decisiones administrativas bajo incertidumbre}
Un parámetro $\theta$ describe una población, como el gasto medio de los clientes de un periodo. Un \textit{estimador} $\hat{\theta}$ es una regla o función de la muestra aleatoria que varía entre muestras; la \textit{estimación} es el valor numérico obtenido al evaluarla en datos observados. El \textit{insesgamiento} exige que la esperanza del estimador coincida con el parámetro, $E(\hat{\theta})=\theta$; el \textit{error cuadrático medio} suma varianza y sesgo al cuadrado, $\operatorname{ECM}(\hat{\theta})=\operatorname{Var}(\hat{\theta})+[B(\hat{\theta})]^2$ \citep[sec.~8.2]{wackerly2010}. La \textit{eficiencia relativa} compara varianzas de estimadores insesgados, y la \textit{consistencia} exige que la probabilidad de exceder cualquier tolerancia fija tienda a cero al crecer $n$ \citep[secs.~9.2--9.3]{wackerly2010}.

La \textit{estimación puntual} resume el parámetro en una sola cifra pero no informa su precisión; la \textit{estimación por intervalos} añade límites calculados con el error estándar y un nivel de confianza $1-\alpha$. El \textit{nivel de confianza} describe la proporción de intervalos que contendrían el parámetro en repeticiones del muestreo bajo sus supuestos; el \textit{margen de error} mide la semiamplitud de un intervalo simétrico, crece con mayor confianza o dispersión y disminuye al aumentar $n$ \citep{holmesBusiness2023,nistHandbook}.

\begin{fichainvestigacion}{Ficha 1. Propiedades del estimador en un caso de compras}
\datoscasocorto{\textbf{Caso y fuentes:} Gasto medio por pedido en compras simuladas (\$180, \$200, \$220 y \$240) y ampliación de 100 a 400 registros \citep[secs.~8.2, 9.2 y 9.3]{wackerly2010}.}
\textbf{Estimador y propiedades.} Promediar es el estimador $\bar{X}$ y \$210 es la estimación puntual, que aproxima pero no iguala necesariamente la media poblacional $\mu$. Con compras independientes de igual distribución, media $\mu$ y varianza finita $\sigma^2$, tanto $\bar{X}_4$ como la primera compra $X_1$ son insesgados. La varianza de $\bar{X}_4$ es $\sigma^2/4$ frente a $\sigma^2$ en $X_1$: el promedio es cuatro veces más eficiente ($\operatorname{ER}=4$) porque aprovecha las cuatro observaciones. Por la desigualdad de Chebyshev, $\bar{X}_n$ es consistente; pasar de 100 a 400 compras reduce el error estándar $\sigma/\sqrt{n}$ a la mitad, mientras que usar solo $X_1$ conserva varianza $\sigma^2$ y no converge al parámetro.
\end{fichainvestigacion}

\begin{fichainvestigacion}{Ficha 2. Estimación puntual e intervalar en áreas funcionales}
\datoscasocorto{\textbf{Ámbito y fuentes:} Finanzas, recursos humanos y mercadotecnia \citep{holmesBusiness2023,bradley2021}.}
\textbf{Aplicación comparada.} En \textit{finanzas}, el gasto medio estimado de \$210 sirve como escenario central de presupuesto maestro, mientras que el intervalo dimensiona reservas de liquidez y sensibilidad ante variaciones de ingresos. En \textit{recursos humanos}, la proporción puntual $\hat{p}$ de empleados que requieren capacitación orienta la programación inicial y su intervalo define el rango de instructores y aulas, sin atribuir esa proporción a cada trabajador ni omitir diferencias de respuesta entre turnos. En \textit{mercadotecnia}, la tasa puntual de conversión resume una campaña piloto y su intervalo evita escalar gastos publicitarios cuando la mejora no supera la variación muestral o excluye canales de venta.
\end{fichainvestigacion}

\begin{fichainvestigacion}{Ficha 3. Caso aplicado: intervalo de ventas semanales [\$18,000, \$22,000]}
\datoscasocorto{\textbf{Planteamiento y fuente:} Media semanal de ventas estimada entre \$18,000 y \$22,000 con 95\% de confianza \citep[sec.~1.3.5.2]{nistHandbook}.}
\textbf{Interpretación gerencial.} El intervalo simétrico tiene centro puntual de \$20,000 y margen de error de \$2,000 ($\pm 10\%$), sin inventar muestra ni desviación no declaradas. Describe valores plausibles para el \textit{promedio poblacional} de ventas semanales $\mu$ bajo el diseño y periodo estudiados; no garantiza que el 95\% de las semanas individuales venda dentro de ese rango ni asigna probabilidad frecuentista a un parámetro fijo ya calculado. Predecir la siguiente semana requeriría un intervalo de predicción y evaluar estacionalidad. Si el punto de equilibrio supera \$18,000, cae dentro del intervalo y la gerencia debe revisar costos, reservas o tamaño muestral antes de comprometer inventario.
\end{fichainvestigacion}

La estimación estadística es fundamental para la gestión administrativa, las finanzas, los recursos humanos y la mercadotecnia porque convierte registros parciales en decisiones con riesgo acotado: permite elegir reglas insesgadas, eficientes y consistentes, separar un escenario central de su margen de error y evitar que una muestra grande pero sesgada se confunda con evidencia concluyente \citep{wackerly2010,nistHandbook,bradley2021}.

\section{Conclusiones}
La estimación aporta soporte administrativo cuando articula población objetivo, propiedades del estimador y lectura rigurosa de la incertidumbre. Insesgamiento, eficiencia y consistencia cumplen funciones distintas y ninguna propiedad algebraica corrige por sí sola un marco incompleto. Sostengo que toda gerencia debe exigir junto con cada cifra puntual su población, supuestos y margen de error, porque decidir solo con un promedio oculta los riesgos trasladados a presupuestos, personal o campañas comerciales. La consecuencia profesional es formular planes con umbrales de revisión y reservas operativas acordes con la amplitud del intervalo estimado.

\renewcommand{\bibsection}{\section*{Referencias}\addcontentsline{toc}{section}{Referencias}}
\setlength{\bibhang}{1.27cm}\setlength{\bibsep}{4pt}
\begin{thebibliography}{4}
\bibitem[Bradley et al.(2021)]{bradley2021}
Bradley, V. C., Kuriwaki, S., Isakov, M., Sejdinovic, D., Meng, X.-L., \& Flaxman, S. (2021). Unrepresentative big surveys significantly overestimated US vaccine uptake. \textit{Nature, 600}(7890), 695--700. \url{https://doi.org/10.1038/s41586-021-04198-4}
\bibitem[Holmes et al.(2023)]{holmesBusiness2023}
Holmes, A., Illowsky, B., \& Dean, S. (2023). \textit{Introductory business statistics} (2.\textsuperscript{a} ed.). OpenStax. \url{https://openstax.org/details/books/introductory-business-statistics-2e}
\bibitem[NIST y SEMATECH(2003)]{nistHandbook}
National Institute of Standards and Technology, \& SEMATECH. (2003). \textit{NIST/SEMATECH e-Handbook of statistical methods}. \url{https://doi.org/10.18434/M32189}
\bibitem[Wackerly et al.(2010)]{wackerly2010}
Wackerly, D. D., Mendenhall, W., \& Scheaffer, R. L. (2010). \textit{Estadística matemática con aplicaciones} (7.\textsuperscript{a} ed.). Cengage Learning Editores.
\end{thebibliography}
\global\itescaactivitycitedfalse



